\section{Aggregate-volume closure recovered}
RECOVERY-STRENGTHEN7 proves the general recursive fixed point matching the
latest concrete ledgers.  If one uniform hereditary all-input phase has local
time exponent $r$, aggregate executed child volume exponent $v$, maximum
child exponent $h<1$, and local-state exponent $p$, then recursive
substitution gives one exact same-machine pair
\[
 C=\max\left\{1,r,\frac{v-h}{1-h}\right\},\qquad
 P=\max\{p,h\}.
\]
The proof uses $\sum_i m_i^C\le(\max_i m_i)^{C-1}\sum_i m_i$ and counts
repeated children repeatedly.  Thus
$(r,v,h,p)=(6/5,11/10,1/2,1/2)$ really would eliminate the inherited
exponent $K$ and yield the joint pair $(6/5,1/2)$.

This strengthens the original count-envelope RECURCLOSE1 theorem, whose
specialization has $v=\eta+h$.  It does not create the required all-input
phase.  A finite partial portfolio must cover every root and actual successor;
RECURCLOSE1, CYCLEGAIN1, STREAMCOMP1 and ORBITBATCH1 account for an already
defined execution graph but do not add transitions at FAIL vertices.

An explicit guarded family remains outside the current favorable marked and
pivot members.  Therefore the recovered fixed point is conditional and the
FLMV endpoint is not reached.  P versus NP remains OPEN.  See
\path{research/agents/proof/2026-10-08_2349_RECOVERY-STRENGTHEN7.tex}.
